\chapter{Introduction}

\texttt{PrimePie} is a Python number theory library with C++ backend, designed for high-performance computation and symbolic experimentation.

It is structured into three primary components:
\begin{enumerate}
    \item Analytic Number Theory
    \item Algebraic Number Theory
    \item Cryptographic Algorithms
\end{enumerate}
The ie project is guided by the following Philosophy:
\begin{itemize}
    \item \textbf{Pythonic interface:} A simple and flexible Python API that feels native and encourages experimentation.
    \item \textbf{High-performance C++ backend:} Core computations are implemented in C++ for speed and scalability.
    \item \textbf{Seamless visualization:} Built-in support for intuitive and customizable visualization of number-theoretic objects and results.
    \item \textbf{Minimal and general abstractions:} Mathematical structures are defined in a general and reusable way, without unnecessary overhead.
\end{itemize}


\section{Philosophy}

\texttt{PrimePie} is designed with several principles in mind:

\paragraph{1. Data Types -- Precision VS. Approximation}
\textcolor{red}{TODO}


\paragraph{2. Parallel Computation -- Scalar Logic, Vector Execution}
All scalar operations in the library are automatically extended to \textbf{vectorized operations}, allowing users to apply the same logic to entire collections of values with minimal effort.

For example, consider the method \texttt{.power(a, n)} which computes $a^n$:

\begin{itemize}
  \item If \texttt{a} is a scalar and \texttt{n} is a scalar, it returns $a^n$.
  \item If \texttt{a} is a list $[a_1, a_2, \dots, a_k]$ and \texttt{n} is a scalar $m$, the function returns $[a_1^m, a_2^m, \dots, a_k^m]$.
  \item If both \texttt{a} and \texttt{n} are lists of the same length, the function performs element-wise exponentiation:
  \[
    \texttt{power(}[a_1, \dots, a_k], [n_1, \dots, n_k]\texttt{)} \rightarrow [a_1^{n_1}, \dots, a_k^{n_k}]
  \]
\end{itemize}
